\section{Coefficient seven and small real residual fields}
\label{sec:coefficient-seven-real-residual}

\paragraph{Original coefficients and the exact finite scope.}
Put $K=\mathbb Q(\omega)$ and $\eta_K=\chi_{-3}$.
Let $E/\mathbb Q_7$ be any finite extension, let $O$ be its original
ring of integers, and put $D_O=E/O$.
Let $c_7$ be the ordinary cyclotomic character and let $\tau_7$ be
its Teichm\"uller lift, so that $\overline{\tau_7}=\overline{c_7}$
and $\tau_7$ has order six. All its values already lie in
$\mathbb Z_7^\times$, hence in this unchanged $O$.
For $d\in\{-4,8,-8\}$, $\chi_d$ denotes the primitive quadratic
Dirichlet character of the indicated discriminant. Set
\[
 \begin{gathered}
 \mathcal D=\{1,\chi_{-4},\chi_8,\chi_{-8}\},\\
 \mathcal A_e=
 \{(\delta,j):\delta\in\mathcal D,\ 0\leq j<6,
       \ \delta\ne\chi_8\ \hbox{or }j-e\hbox{ is odd}\}.
 \end{gathered}
 \tag{R7.1}
\]
For an integer $e\geq1$ and $(\delta,j)\in\mathcal A_e$, let
$k\in\{0,1\}$ be the unique choice for which
\[
 \begin{aligned}
 \psi^{\rm nc}&=\delta\tau_7^j\eta_K^k,
 &\psi^{\rm nc}(-1)&=(-1)^e,\\
 \psi^{\rm cr}&=\psi^{\rm nc}\eta_K,\\
 a^{\rm nc}&=\psi^{\rm nc}c_7^{-e},
 &a^{\rm cr}&=\psi^{\rm cr}c_7^{-e}.
 \end{aligned}\tag{R7.2}
\]
Thus complex conjugation acts by $+1$ and $-1$ on the two full
coefficients, respectively. The terms noncritical and critical
refer to precisely these two parities, as in (CP.3).

For a number field $M$, write $\mathcal S_7(M,A)$ for the kernel
of continuous global $H^1(M,A)$ restricted to full $H^1(M_v,A)$
at every $v\mid7$ and to $H^1(I_v,A)$ at every other finite place.
Include real-place $H^1$ when present; it is zero on a
$7$-primary coefficient. These are the full local conditions,
on the original coefficient. A finite localization containing
all coefficient-ramification primes gives the same kernel.

\begin{theorem}[AC-G197: real residual descent at coefficient seven]
\label{AC-G197}
For every original $O$ above, every $e\geq1$, and every
$(\delta,j)\in\mathcal A_e$,
\[
             \mathcal S_7(\mathbb Q,D_O(a^{\rm nc}))=0.
 \tag{R7.3}
\]
On this same coefficient there is an exact identification
\[
 \mathcal S_7\bigl(K,D_O(\delta\tau_7^j c_7^{-e})\bigr)
       \simeq \mathcal S_7(\mathbb Q,D_O(a^{\rm cr})).
 \tag{R7.4}
\]
In particular (R7.4) does not assert original-$K$ strict zero.

In the actual host and original coefficient of (FC.1)--(FC.2),
assume $h\geq2$ and let $T\geq1$ be an integer. All the triples
\[
 Y^{\rm real}_7(T)=
 \{(7,(\delta\tau_7^j)|_{G_K},e):1\leq e<T,
                                      (\delta,j)\in\mathcal A_e\}
 \tag{R7.5}
\]
are actual trivial-cubic labels, with all undisplayed finite
factors trivial. They are distinct and
\[
                        |Y^{\rm real}_7(T)|=21(T-1).
 \tag{R7.6}
\]
For any subset $Y'\subseteq Y^{\rm real}_7(T)$ and $T\geq2$,
the already proved critical aggregate bound applies once:
\[
 \sum_{(7,\eta,e)\in Y'}
       \frac{\log\#\mathcal S_{7,\eta}(e)}{b_7}
       \leq\mathcal B_{\rm cr}(F,T)\leq C\mathcal P(F,T),
 \tag{R7.7}
\]
with exactly the host, polynomial and constant scope of AC-G191.
The nonnegative remainder of (CP.3) therefore satisfies, for every
retained set $J$ of actual triples,
\[
             \mathcal R_J=
                    \mathcal R_{J\setminus Y^{\rm real}_7(T)}.
 \tag{R7.8}
\]
There is no positive-derived-index splitting or zero assertion.
No claim is made here for $\delta=\chi_8$ with $j-e$ even.
\end{theorem}

\begin{proof}
\proofstep{1. Four exact finite class-quotient certificates.}
Put $R=\mathbb Q(\zeta_7)^+$ and $B=\mathbb Q(\zeta_{84})^+$.
The following four totally real fields are the only class inputs:
\[
                         B,\quad R\sqrt6,\quad R\sqrt{14},
                                      \quad R\sqrt{42}.
 \tag{R7.9}
\]
Here $R\sqrt d$ means the compositum $R\,\mathbb Q(\sqrt d)$.
Their explicit defining polynomials, in the displayed order, are
\[
 \begin{aligned}
 P_B(X)={}&X^{12}-11X^{10}+44X^8-78X^6+60X^4-16X^2+1,\\
 P_6(X)={}&X^6+2X^5-21X^4-30X^3+98X^2+64X-71,\\
 P_{14}(X)={}&X^6+2X^5-45X^4-62X^3+562X^2+368X-1847,\\
 P_{42}(X)={}&X^6+2X^5-129X^4-174X^3+5210X^2+3448X-65519.
 \end{aligned}\tag{R7.10}
\]
The exact identity
$\Phi_{84}(Z)=Z^{12}P_B(Z+Z^{-1})$ identifies the first field.
For a primitive $84$th root $z$, the real number $z+z^{-1}$
has degree $\varphi(84)/2=12$: the nonreal $z$ is quadratic
over its real trace field. This proves the asserted field and degree.

Likewise $P_R(Y)=Y^3+Y^2-2Y-1$ follows from
$\Phi_7(Z)=Z^3P_R(Z+Z^{-1})$.
For $d=6,14,42$ the polynomial in (R7.10) is exactly
\[
                P_d(X)=\operatorname{Res}_Y
                          (P_R(Y),(X-Y)^2-d).
 \tag{R7.11}
\]
The cubic field $R$ and the quadratic field $\mathbb Q(\sqrt d)$
are disjoint. If $\alpha=\zeta_7+\zeta_7^{-1}$, a Galois
automorphism fixing $\alpha+\sqrt d$ would satisfy
$\sigma(\alpha)-\alpha=(1-\varepsilon)\sqrt d$.
For $\varepsilon=-1$ this contradicts disjointness, and for
$\varepsilon=1$ it forces $\sigma=1$. Thus this element has
degree six and (R7.11) defines precisely the claimed compositum.
The saved exact polynomial script independently checks these identities.

For each polynomial a fresh \texttt{bnfinit(P,1)} was followed by
\texttt{bnfcertify(B,1)}. The exact printed outputs are
\[
 \begin{array}{c|c|c|c|c}
 \text{field}&\text{degree}&\text{signature}&\text{computed quotient}
                                      &\text{flag one}\\\hline
 B&12&(12,0)&1&1\\
 R\sqrt6&6&(6,0)&1&1\\
 R\sqrt{14}&6&(6,0)&1&1\\
 R\sqrt{42}&6&(6,0)&C_2&1
 \end{array}\tag{R7.12}
\]
The applied PARI manual states that flag one certifies that the
actual class group is a quotient of the computed one.
It does not certify a full unit basis. Consequently all four
actual class groups have trivial $7$-primary part, unconditionally.
No candidate class number is substituted for this successful
quotient certificate, and no general GRH assumption is used.
The scripts, complete outputs and binary field data are retained.
An exploratory degree-$24$ real-$168$ computation is not an input
to this argument; its pending certification supplies no conclusion.

\proofstep{2. The actual even residual field and its class group.}
Let $r=j-e\pmod6$, write $\epsilon(\delta)=0$ or $1$ according
as $\delta(-1)=1$ or $-1$, and take
$k\equiv\epsilon(\delta)+r\pmod2$.
Define the finite character
\[
       \rho=\delta\tau_7^r\eta_K^k,
       \qquad M_\rho=\overline{\mathbb Q}^{\ker\rho}.
 \tag{R7.13}
\]
Its reduction is exactly $\overline{a^{\rm nc}}$.
Its order divides six; reduction of these roots of unity at
seven is injective, so this is the literal residual splitting field.
It is real because $\rho(-1)=1$, cyclic over $\mathbb Q$, and
$[M_\rho:\mathbb Q]\mid6$.

Here are the exact class covers for all admitted cases:
\[
 \begin{array}{c|c|c}
 \delta& r\text{ condition}&M_\rho\text{ is contained in}\\\hline
 1\text{ or }\chi_{-4}&\text{any}&B\\
 \chi_{-8}&r\text{ even}&R\sqrt6\\
 \chi_{-8}&r\text{ odd}&R\sqrt{14}\\
 \chi_8&r\text{ odd}&R\sqrt{42}.
 \end{array}\tag{R7.14}
\]
For the first row, $\rho$ is even and has conductor dividing $84$.
For the other rows use $\tau_7^3=\chi_{-7}$ and
\[
 \chi_{-8}\eta_K=\chi_{24},\qquad
 \chi_{-8}\chi_{-7}=\chi_{56},\qquad
 \chi_8\eta_K\chi_{-7}=\chi_{168}.
 \tag{R7.15}
\]
The quadratic fields of these three characters are respectively
$\mathbb Q(\sqrt6)$, $\mathbb Q(\sqrt{14})$ and
$\mathbb Q(\sqrt{42})$. The remaining even power of $\tau_7$
factors through the cubic field $R$.
These identities can also be checked directly from the quadratic
Dirichlet values; no field with a different ramification set is substituted.

Each cover in (R7.14) has degree six or twelve, prime to seven.
For an extension $V/M_\rho$, norm composed with ideal extension
on class groups is multiplication by $[V:M_\rho]$.
Its restriction to the $7$-primary subgroup is injective.
The successful class-quotient certificates in Step 1 therefore give
\[
                         \operatorname{Cl}(M_\rho)[7]=0.
 \tag{R7.16}
\]
This is a statement about the smaller residual field, and not
an inference that a larger class cover has a unique prime above seven.

\proofstep{3. The unique prime is checked on that residual field.}
All fields in (R7.13) lie in $L=\mathbb Q(\zeta_{168})$.
The prime-to-seven factor $\mathbb Q(\zeta_{24})$ is unramified
at seven, with Frobenius $7\bmod24$ of order two.
The factor $\mathbb Q(\zeta_7)$ is totally ramified of degree six.
Thus a decomposition group in $\operatorname{Gal}(L/\mathbb Q)$
is the product of $\langle7\bmod24\rangle$ and the full
modulo-seven factor. Its image under $\rho$ is generated by
\[
          \tau_7^r(I_7)\quad\hbox{and}\quad\delta(7),
                    \qquad \eta_K(7)=1.
 \tag{R7.17}
\]
Write $d_r=6/\gcd(6,r)$, with $d_0=1$.
The first group has order $d_r$.
For $\delta=1$, the complete image of $\rho$ has order $d_r$:
when $r$ is odd, the sign from $\eta_K$ is already in this
even-order image; when $r$ is even, $k=0$.
Hence inertia alone gives the full image.

For each nontrivial $\delta\in\mathcal D$, independent
prime-to-seven conductor factors show that the complete image
has order $\operatorname{lcm}(2,d_r)$.
The exact values are
\[
                 \chi_{-4}(7)=\chi_{-8}(7)=-1,
                         \chi_8(7)=1.
 \tag{R7.18}
\]
For $\delta=\chi_{-4},\chi_{-8}$, (R7.17) is therefore the
entire image for every $r$. For $\delta=\chi_8$ it is the
entire image exactly when $d_r$ is even, equivalently $r$ is odd.
These are precisely the cases in (R7.1).
The decomposition group in the abelian field $M_\rho$ is thus
its entire Galois group, proving that $M_\rho$ has one prime above seven.

For comparison, $L^+$ itself has two such primes: the full
decomposition group has order twelve and does not contain
complex conjugation, while $[L^+:\mathbb Q]=24$.
We have not applied a one-prime lemma to $L^+$.
In the excluded $\chi_8$ cases, the residual decomposition image
has index two. This identifies the limitation of this argument;
it is not a nonvanishing assertion for the excluded strict groups.

\proofstep{4. One-prime local generation with the full local witness.}
We repeat the odd-prime lemma used in AC-G193 at the exact fields
needed here. Let $M$ have one prime $v$ above an odd prime $p$
and let $\operatorname{Cl}(M)[p]=0$.
Let $\Gamma$ be the Galois group of the maximal extension
unramified at finite primes other than $v$, and let
$P=\Gamma^{(p)}$ be its maximal pro-$p$ quotient.
The compact image of $G_{M_v}$ in $P$ is closed.
If proper, it is proper in a finite $p$-group quotient of $P$.
Every proper subgroup of a finite $p$-group lies in a normal
subgroup of index $p$. One proof takes a maximal subgroup and
uses the nontrivial center inductively: if the center is not
contained in it, their product is the group; otherwise quotient
by the center. The resulting $C_p$ quotient gives a nontrivial
cyclic extension of $M$ unramified away $v$ and split at $v$.
It is unramified at all finite primes, and split at real primes
because a real decomposition group has order at most two.
Ordinary Hilbert class-field reciprocity contradicts
$\operatorname{Cl}(M)[p]=0$. The local image is therefore all of $P$.

Let $a:G_M\to1+\mathfrak m_O$ be continuous and unramified
at finite primes other than $v$, using any original $p$-adic
coefficient ring $O$. If $z:G_M\to(E/O)(a)$ is a strict cocycle,
its image is finite. Full local zero supplies a witness $x\in E/O$
such that $z(g)=(a(g)-1)x$ on $G_{M_v}$.
Choose a power $\pi^m$ of an original uniformizer killing
both $x$ and the entire image of $z$.
The scalar action on $(E/O)[\pi^m]$ has finite $p$-group image $U_m$.
The cocycle homomorphism
\[
 g\longmapsto\bigl(z(g),a(g)|_{(E/O)[\pi^m]}\bigr)
          \ \in (E/O)[\pi^m]\rtimes U_m
 \tag{R7.19}
\]
has a finite $p$-group image. Away from $v$ inertia acts trivially
on the coefficient, so zero inertia cohomology means the cocycle
itself vanishes there. Hence (R7.19) factors through $\Gamma$
and then $P$. Its local and global images coincide.
For every global element choose a local element with the same
image; their actions on the chosen witness agree.
The local coboundary formula is consequently global, proving
\[
                      \mathcal S_p(M,(E/O)(a))=0.
 \tag{R7.20}
\]
This argument retains the actual finite local witness; it does
not infer a common coboundary merely from generation by several
different local subgroups.

\proofstep{5. Descent on the unchanged coefficient.}
For $M=M_\rho$, reduction of $a^{\rm nc}$ is trivial by its
literal definition in (R7.13). Thus its full restriction has
values in $1+\mathfrak m_O$ without any coefficient extension.
It is unramified outside seven on $G_M$.
Indeed $c_7$ is unramified at every other finite prime.
On its inertia subgroup the finite character $\psi^{\rm nc}$
has order dividing six and reduction one.
Reduction is injective on roots of unity of order prime to seven:
if $z^n=1$, $z\equiv1\pmod{\mathfrak m_O}$ and $7\nmid n$,
then $1+z+\cdots+z^{n-1}$ is a unit, forcing $z=1$.
This proves the claimed full, not just residual, unramifiedness.

Steps 2--3 give the class and unique-prime hypotheses of Step 4.
Therefore the restricted strict group over $M$ is zero.
A rational strict class restricts into it: full zero at seven
restricts to full zero at the actual prime above seven, and
zero on every other inertia group restricts to zero on each
inertia subgroup. Restriction followed by corestriction on
$H^1$ is multiplication by $[M:\mathbb Q]\mid6$.
This is invertible on a $7$-primary group. Restriction is
injective, proving (R7.3) for the original $O$, however ramified.

\proofstep{6. Both rational extensions and every full local condition.}
The two finite characters in (R7.2) restrict to the same character
$\delta\tau_7^j$ of $G_K$.
The regular $C_2$ factor of the induced original lattice splits
by $(1\pm c)/2$, since $2\in O^\times$. Thus
\[
 \operatorname{Ind}_{G_K}^{G_{\mathbb Q}}
            O(\delta\tau_7^j c_7^{-e})
       \simeq O(a^{\rm nc})\oplus O(a^{\rm cr}).
 \tag{R7.21}
\]
Its change-of-basis determinant is a unit times two.
Tensoring with $E/O$ and applying Shapiro gives a direct sum
on global cohomology, without a lattice-index error.

At seven, local Shapiro retains both actual primes of $K$, since
seven splits in $K$, and compares full local $H^1$ groups.
At three, $K/\mathbb Q$ is ramified quadratic; its inertia
quotient is the full $C_2$. The inertia Mackey decomposition
therefore has one double coset, with intersection the actual
$K$-inertia subgroup. Shapiro for this subgroup identifies the
zero inertia conditions on the two sides of (R7.21).
At every other finite prime the quadratic extension is unramified;
its two or conjugate inertia summands give exactly the conditions
at all primes of $K$ above it. Real-place $H^1$ vanishes at seven.
Consequently the full strict group over $K$ is the direct sum
of the two rational strict groups. Removing the zero summand
from (R7.3) proves (R7.4).

If the host localization adds $q,d,269$ or other primes, the same
comparison remains literal. At a prime where the coefficient is
unramified, zero inertia means that the cocycle vanishes on inertia;
its finite cocycle homomorphism kills the closed normal subgroup
generated by these inertia groups. Inflation therefore identifies
the strict kernels before and after that enlargement.
This is also the finite-\'etale/Galois-cohomology convention
retained in (FC.2). No completion or coefficient normalization changes.

\proofstep{7. Actual labels, counts and the one critical charge.}
The actual host of (FG.3) contains $\mathbb Q(\zeta_{2^{h+2}})$
and $\mathbb Q(\zeta_7)$. For $h\geq2$ it therefore contains
all four quadratic characters in $\mathcal D$ and all six
modulo-seven characters. Their values are already in the
specified original $O_7$, whose root-of-unity exponent is
$3\cdot2^h$. The cubic factor and the $q,d$ factors in (R7.5)
are trivial. Ratios of two displayed rational labels have
conductor dividing $56$. Such a ratio cannot equal $\eta_K$,
which is ramified at three; hence restricting to $G_K$ does
not identify distinct displayed labels.

At each weight the three characters $1,\chi_{-4},\chi_{-8}$
allow all six values of $j$, while $\chi_8$ allows exactly
the three values with $j-e$ odd. This proves (R7.6), including
the empty case $T=1$. The coefficient-seven residual-identity
label $\tau_7^e$ from AC-G195 occurs once in this set at every
weight. We do not make an unproved disjointness assertion about
any larger previously removed union here.

For each admitted actual triple, (R7.4) identifies its full
strict group with its unique critical rational component on
the same original $O_7$. Their normalized logarithmic sizes
use the same $b_7$. An arbitrary subset is bounded by the
single full nonnegative critical aggregate in AC-G191, proving
(R7.7) with no factor twenty-one.
The corresponding noncritical summands of (CP.3) are zero
term by term by (R7.3), giving (R7.8) for every retained set.
This equality does not separate positive parts of derived
indices, and it is independent of a choice of common window.
The later actual-window comparison must retain its one total
local/torsion error and the existing one critical aggregate.
\end{proof}

\paragraph{Source and computation scope.}
The four finite class assertions use only the successful
PARI flag-one quotient certificates and the complete applied
manual paragraph. The one-prime proof is supplied above from
the ordinary Hilbert class-field statement, the maximal
$S$-ramified extension construction, and elementary finite
$p$-group and cocycle arguments. Those primary source scopes
and the exact AC-G193/AC-G195 comparisons are retained in the
source manifest; no new proof of global reciprocity is claimed.
The original-host and critical-cost interfaces are exactly
(FG.3), (FC.1)--(FC.2) and AC-G191. Exploratory real-$168$
candidate output is preserved separately and unused.
